Transient simulations will apply the improved \gls{DNP} group parameters in the \gls{MSRE} to demonstrate how they affect power and temperature responses.
They are critical for validating the \gls{DNP} group parameters using \gls{MSRE} transients and for investigating transients affecting safety, security, and safeguards design decisions.
Section \ref{sec:transimpact} discusses the proposed simulations in detail.
These simulations will use OpenMC to generate multigroup cross sections at three temperatures: a higher than nominal, a lower than nominal, and the nominal \gls{MSRE} temperature.

For safety design, the \gls{ULOF}, \gls{UTOP}, \gls{ULOHS}, and \gls{FSOC} will be investigated for various reactor parameters to assess the impact of the improved \gls{DNP} group parameters.
For security design, sabotage-initiated transients will build on those studied for safety design.
For safeguards design, comparisons are made between diversion scenarios.
The diversion scenarios alter the fuel composition, affecting the effective delayed neutron fraction and reactor behavior during transients \cite{wheeler2021signatures}.
The improved \gls{DNP} group parameters will be used to investigate the extent of \gls{SNM} removal which could remain undetected while using previous \gls{DNP} group parameters.

The kinetics simulations will use Moltres models of the \gls{MSRE} to evaluate how the group parameters affect reactivity and temperature responses.
These results can be compared with other analyses from the literature \cite{ho2023exploring, singh2018nonlinear}.
The results will identify which transients are more sensitive to the improved data and highlight potential concerns.
%A \gls{PRK} model can also be used for simple reactivity insertion tests with uncertainty tracking to more straightforwardly demonstrate differences between the \gls{DNP} group parameters.
